% GENERATED FILE. Edit canonical lessons, then run npm run generate:books.
% canonical-source-hash: fnv1a64:24528a79ea7b785e
% canonical-lessons: GE-C348-verantwortung, GE-C348-verantwortlich, GE-C348-rechtfertigen, GE-C348-ausrede, GE-C348-deswegen

\chapter{Wer ist verantwortlich?}
\label{ch:ge-wer-ist-verantwortlich}
\hlchaptermodality{348}

\begin{chapteropening}
\textbf{By the end of this chapter:} \emph{I can say who is responsible, justify what I did, and recognize an excuse.}

Complete the last lesson of chapter 348: I can say who is responsible, justify what I did, and recognize an excuse.
\end{chapteropening}

\section[die Verantwortung]{die Verantwortung — the responsibility}

\label{lesson:GE-C348-verantwortung}



\begin{quote}
Before the new one: say the German for to avoid, then the German for why.
\end{quote}

\subsection*{You'll want to know: die Verantwortung}
\textbf{die Verantwortung} — \textquotedblleft{}the responsibility\textquotedblright{}.

\begin{quote}
Die Eltern tragen die Verantwortung für ihre Kinder.
\end{quote}

\begin{quote}
\emph{Parents bear responsibility for their children.}
\end{quote}

\emph{Die Verantwortung übernehmen} — to take responsibility.

\begin{cousinweb}
\emph{Verantwortung} is the noun of \emph{verantworten} \textquotedblleft{}to answer for\textquotedblright{}, built on \emph{Antwort} \textquotedblleft{}answer\textquotedblright{}, from Old High German \emph{antwurti} and Proto-Germanic \emph{andawurdiją} (EtymDB), which is built on \emph{Wort} \textquotedblleft{}word\textquotedblright{}.
\end{cousinweb}

\begin{grammarlens}[title={What you've built}]
One more word for telling, explaining and describing at B1.
\end{grammarlens}

\subsection*{Your turn}
\begin{itemize}
\raggedright
  \item \emph{Say it:} \emph{die Verantwortung}
  \item \emph{Say it:} \emph{die Verantwortung}, once more
  \item \emph{Say it:} a sentence of your own with \emph{die Verantwortung}
  \item \emph{Recall:} say the German for to avoid, then the German for why, then say \emph{die Verantwortung} again
\end{itemize}

\begin{tcolorbox}[breakable,colback=teal!4,colframe=teal!35!black,title={Before you move on}]
What does die Verantwortung mean? (\textquotedblleft{}The responsibility\textquotedblright{}.) Say it once more.
\end{tcolorbox}
\section[verantwortlich]{verantwortlich — responsible}

\label{lesson:GE-C348-verantwortlich}



\begin{quote}
Before the new one: say the German for why, then the German for the responsibility.
\end{quote}

\subsection*{You'll want to know: verantwortlich}
\textbf{verantwortlich} — \textquotedblleft{}responsible\textquotedblright{}.

\begin{quote}
Mein Bruder ist heute für die Küche verantwortlich.
\end{quote}

\begin{quote}
\emph{My brother is responsible for the kitchen today.}
\end{quote}

\emph{Verantwortlich für} + accusative.

\begin{cousinweb}
\emph{verantwortlich} is built with \emph{-lich} on \emph{verantworten} \textquotedblleft{}to answer for\textquotedblright{}, from \emph{Antwort} \textquotedblleft{}answer\textquotedblright{}, Old High German \emph{antwurti} (EtymDB): someone who must answer for a thing.
\end{cousinweb}

\begin{grammarlens}[title={What you've built}]
One more word for telling, explaining and describing at B1.
\end{grammarlens}

\subsection*{Your turn}
\begin{itemize}
\raggedright
  \item \emph{Say it:} \emph{verantwortlich}
  \item \emph{Say it:} \emph{verantwortlich}, once more
  \item \emph{Say it:} a sentence of your own with \emph{verantwortlich}
  \item \emph{Recall:} say the German for why, then the German for the responsibility, then say \emph{verantwortlich} again
\end{itemize}

\begin{tcolorbox}[breakable,colback=teal!4,colframe=teal!35!black,title={Before you move on}]
What does verantwortlich mean? (\textquotedblleft{}Responsible\textquotedblright{}.) Say it once more.
\end{tcolorbox}
\section[rechtfertigen]{rechtfertigen — to justify (oneself)}

\label{lesson:GE-C348-rechtfertigen}



\begin{quote}
Before the new one: say the German for the responsibility, then the German for responsible.
\end{quote}

\subsection*{You'll want to know: rechtfertigen}
\textbf{rechtfertigen} — \textquotedblleft{}to justify (oneself)\textquotedblright{}.

\begin{quote}
Er musste sich vor seinem Chef rechtfertigen.
\end{quote}

\begin{quote}
\emph{He had to justify himself to his boss.}
\end{quote}

Inseparable, although it looks like a compound: \emph{er rechtfertigt sich, er hat sich gerechtfertigt}.

\begin{cousinweb}
\emph{rechtfertigen} is built from \emph{recht} \textquotedblleft{}right\textquotedblright{}, from Proto-Germanic \emph{rehtaz} (EtymDB), and \emph{fertigen} \textquotedblleft{}to make ready\textquotedblright{}, from \emph{fertig}: to set something right.
\end{cousinweb}

\begin{grammarlens}[title={What you've built}]
One more word for telling, explaining and describing at B1.
\end{grammarlens}

\subsection*{Your turn}
\begin{itemize}
\raggedright
  \item \emph{Say it:} \emph{rechtfertigen}
  \item \emph{Say it:} \emph{rechtfertigen}, once more
  \item \emph{Say it:} a sentence of your own with \emph{rechtfertigen}
  \item \emph{Recall:} say the German for the responsibility, then the German for responsible, then say \emph{rechtfertigen} again
\end{itemize}

\begin{tcolorbox}[breakable,colback=teal!4,colframe=teal!35!black,title={Before you move on}]
What does rechtfertigen mean? (\textquotedblleft{}To justify (oneself)\textquotedblright{}.) Say it once more.
\end{tcolorbox}
\section[die Ausrede]{die Ausrede — the excuse}

\label{lesson:GE-C348-ausrede}



\begin{quote}
Before the new one: say the German for responsible, then the German for to justify.
\end{quote}

\subsection*{You'll want to know: die Ausrede}
\textbf{die Ausrede} — \textquotedblleft{}the excuse\textquotedblright{}.

\begin{quote}
Das ist doch nur eine Ausrede!
\end{quote}

\begin{quote}
\emph{That is only an excuse!}
\end{quote}

Plural: \emph{die Ausreden}. \emph{Eine Ausrede finden} — to find an excuse.

\begin{cousinweb}
\emph{Ausrede} is built from \emph{aus} \textquotedblleft{}out\textquotedblright{} and \emph{Rede} \textquotedblleft{}speech\textquotedblright{}, from the family of \emph{reden} \textquotedblleft{}to talk\textquotedblright{}, Old High German \emph{redōn} (EtymDB): talk that gets you out of trouble.
\end{cousinweb}

\begin{grammarlens}[title={What you've built}]
One more word for telling, explaining and describing at B1.
\end{grammarlens}

\subsection*{Your turn}
\begin{itemize}
\raggedright
  \item \emph{Say it:} \emph{die Ausrede}
  \item \emph{Say it:} \emph{die Ausrede}, once more
  \item \emph{Say it:} a sentence of your own with \emph{die Ausrede}
  \item \emph{Recall:} say the German for responsible, then the German for to justify, then say \emph{die Ausrede} again
\end{itemize}

\begin{tcolorbox}[breakable,colback=teal!4,colframe=teal!35!black,title={Before you move on}]
What does die Ausrede mean? (\textquotedblleft{}The excuse\textquotedblright{}.) Say it once more.
\end{tcolorbox}
\section[deswegen]{deswegen — that is why, for that reason}

\label{lesson:GE-C348-deswegen}



\begin{quote}
Before the new one: say the German for to justify, then the German for the excuse.
\end{quote}

\subsection*{You'll want to know: deswegen}
\textbf{deswegen} — \textquotedblleft{}that is why, for that reason\textquotedblright{}.

\begin{quote}
Ich habe verschlafen. Deswegen bin ich zu spät gekommen.
\end{quote}

\begin{quote}
\emph{I overslept. That is why I was late.}
\end{quote}

The verb comes right after it, as after \emph{deshalb}: \emph{deswegen bin ich …}.

\begin{cousinweb}
\emph{deswegen} joins \emph{des}, the old genitive \textquotedblleft{}of that\textquotedblright{}, to \emph{wegen} \textquotedblleft{}because of\textquotedblright{}, which grew from \emph{von … wegen} \textquotedblleft{}on the part of\textquotedblright{}. EtymDB has the word but no history for it.
\end{cousinweb}

\begin{grammarlens}[title={What you've built}]
One more word for telling, explaining and describing at B1.
\end{grammarlens}

\subsection*{Your turn}
\begin{itemize}
\raggedright
  \item \emph{Say it:} \emph{deswegen}
  \item \emph{Say it:} \emph{deswegen}, once more
  \item \emph{Say it:} a sentence of your own with \emph{deswegen}
  \item \emph{Recall:} say the German for to justify, then the German for the excuse, then say \emph{deswegen} again
\end{itemize}

\begin{tcolorbox}[breakable,colback=teal!4,colframe=teal!35!black,title={Before you move on}]
What does deswegen mean? (\textquotedblleft{}That is why, for that reason\textquotedblright{}.) Say it once more.
\end{tcolorbox}
